With the synaptic drive held fixed, both $r$ and $v$ relax, over a population time
$\tau_{\rm rel}$, to the equilibrium of Eq.~\eqref{eq:mpr}, so $r\simeq F(Q)/\tau_m$.
If $\tau_{\rm rel}$ is much shorter than the synaptic time $\tau_s$
(Eq.~\eqref{eq:slow_synapse}, Appendix~\ref{app:models}), the fast variables can be
eliminated at every instant while the synapses keep their dynamics, and NMM2 reduces to
the static curve with dynamically filtered recurrence
\cite{ruffini2021nmm2,clusella2023}. Applying the static curve to carriers requires the separate condition
$\omega_c\tau_{\rm rel}\ll1$; one population can track synaptic input quasistatically
while responding dynamically to kHz input.
Recurrent excitatory--inhibitory circuits resonate without intrinsic
cellular resonance \cite{brunelhakim1999,brunelwang2003}. Such synaptic
modes survive the fast-membrane limit, whereas modes requiring finite
membrane dynamics need not.

\paragraph{Field coupling and recurrent gain in the slow limit.}
The field input survives this reduction. For one self-coupled population,
$I=I_b+\tau_mJu+\lambda_I E$, where $I_b$ is the normalized background input, $u$ the
postsynaptic potential divided by its steady gain, and $J$ the dimensionless
self-coupling. For a fixed field direction, $\lambda_I=\kappa_E/I_s$, with $\kappa_E$
the equivalent current per unit field and $I_s$ the current unit; a cellular
voltage-polarization coefficient requires an impedance conversion before entering this
current-driven model. At equilibrium $u=r$, and a field change $\delta E$ acts directly and through the
recurrent rate,
\begin{equation}
\delta r=\frac{F'(Q_*)}{\tau_m}\,\big(\lambda_I\,\delta E+\tau_mJ\,\delta r\big),
\end{equation}
and solving for $\delta r$ gives
\begin{equation}
\frac{\partial r_*}{\partial E}
=\frac{\lambda_I F'(Q_*)}{\tau_m[1-JF'(Q_*)]}.
\label{eq:slow-field-gain}
\end{equation}
This stationary derivative does not require slow synapses; it compares the closed loop
with the clamped response $\lambda_IF'(Q_*)/\tau_m$, with $u$ held fixed, at the same
$Q_*$. On a stable equilibrium, positive self-feedback amplifies the steady response by
$1/[1-JF'(Q_*)]$ and negative feedback suppresses it; $JF'(Q_*)<1$ is necessary there but,
with finite membrane time, not a complete stability test. The bare curve $F$ takes total drive as its
argument. A transfer fitted from external drive to settled rate already includes this
feedback and must not be coupled through the same loop again. A JR-like
description retains this gain when input conversion, local slope, and feedback are
matched \cite{ruffini2021nmm2,clusella2023}; its connections between populations do not
automatically specify the same within-population self-coupling.

